\renewcommand{\arraystretch}{1.4}
\begin{longtable}{ >{\raggedright\arraybackslash}p{3.5cm} >{\raggedright\arraybackslash}p{7.5cm} >{\raggedright\arraybackslash}p{3.5cm}}


    \hline
    \textbf{Termine} & \textbf{Descrizione} & \textbf{Termini correlati} \\
    \hline
    \endfirsthead

    \hline
    \textbf{Termine} & \textbf{Descrizione} & \textbf{Termini correlati} \\
    \hline
    \endhead

    \hline
    \multicolumn{3}{r}{{Continua nella pagina successiva}} \\
    \endfoot

    \hline
    \caption{Glossario} \label{tab:glossario}
    \endlastfoot

    \textbf{Insegnamento} & Corso universitario identificato da un codice e caratterizzato da crediti, descrizione e prerequisiti. & Edizione, Prerequisito, Piano di studio \\

    \textbf{Edizione} & Erogazione di un insegnamento in un determinato anno accademico e periodo didattico. Può essere tenuta da uno o più docenti. & Insegnamento, Docente, Lezione \\

    \textbf{Lezione} & Singolo appuntamento all'interno del calendario settimanale di un'edizione. È identificata univocamente dalla terna giorno, fascia oraria e aula. & Edizione \\

    \textbf{Prerequisito} & Insegnamento da superare prima di poter sostenere l'esame di un altro insegnamento. & Insegnamento \\

    \textbf{Studente} & Persona iscritta all'università, identificata dal numero di matricola, associata a un unico corso di laurea. & Piano di studio, Esame, Corso di Laurea \\

    \textbf{Piano di studio} & Insieme dei corsi scelti dallo studente per il proprio percorso accademico. & Studente, Insegnamento \\

    \textbf{Esame} & Tentativo sostenuto da uno studente per un determinato insegnamento, costituito da un'unica prova d'esame. Memorizza il punteggio ottenuto, inclusi i voti insufficienti (< 18). & Studente, Insegnamento \\

    \textbf{Docente} & Figura accademica afferente a un dipartimento, abilitata all'insegnamento. & Edizione, Insegnamento, Dipartimento \\

    \textbf{Corso di Laurea} & Percorso di studi specifico frequentato dallo studente (es. Informatica triennale). & Studente \\

    \textbf{Dipartimento} & Struttura accademica di afferenza per i docenti (es. Dipartimento di Scienze Matematiche, Informatiche e Fisiche). & Docente \\

\end{longtable}